\section{Discussion}\label{ch:discussion}

The experimental validation of SolEvolve across diverse combinatorial domains demonstrates that closed-loop semantic feedback improves LLM-driven mathematical discovery. Unlike traditional workflows where the user bears the burden of verification and guidance, our findings show that an autonomous feedback mechanism can navigate rugged search landscapes more effectively than fitness-only approaches (Section~\ref{sec:ablation}). The semantic analysis addresses a known limitation of LLMs---their propensity for plausible but incorrect reasoning---by providing structured, interpretable feedback.

Our experiments across three combinatorial domains---binary linear codes, self-orthogonal embeddings, and generalized Lucas cubes---demonstrate that this architecture generalizes beyond single-problem optimization. In the domain of error-correcting codes, the system rediscovered benchmark parameters like $[32,14,8]$ and identified $[43,10,16]$ codes with improved weight distributions. The discovery of 3-inequivalent binary $[16,6,6]$ codes illustrates the framework's ability to explore structural diversity in solution spaces.

The extension to self-orthogonal embeddings and generalized Lucas cubes demonstrates the framework's applicability across different mathematical structures. By achieving the Gram-rank lower bound for embeddings over $\mathbb{F}_2, \mathbb{F}_3, \mathbb{F}_4,$ and $\mathbb{F}_5$, SolEvolve provides constructive proofs of optimality. Table~\ref{tab:so-results} summarizes these discoveries, detailing the original code, the minimum scalar additions $t$, the resulting embedding parameters, and the proximity to Grassl bounds. Notably, while the baseline rank-one algorithm yielded a binary embedding with distance $d=4$, the SAT strategy incorporated explicit distance constraints to achieve $d=6$ for the binary $[22,11]$ code. Furthermore, SolEvolve generalized this to $\mathbb{F}_3$, extending a ternary BCH $[13,7,5]$ code to a binary $[20,7,9]$ code self-orthogonal code exactly at the Gram-rank theoretical limit ($t=7$).

\begin{table}[t]
\centering
\caption{SAT-based self-orthogonal embeddings across $\mathbb{F}_2,\mathbb{F}_3,\mathbb{F}_4,\mathbb{F}_5$}
\label{tab:so-results}
\resizebox{\columnwidth}{!}{%
\begin{tabular}{lccccccc}
\toprule
Field & Start Code & Inner Product & Gram Rank & Min $t$ & Output $[n',k]$ & $d_{\min}$ & Grassl \\ \midrule
$\mathbb{F}_2$ & $\mathcal{H}_4\,[15,11,3]$ & Dot & 7 & 7 & $[22,11]$ & 6 & 7 \\
$\mathbb{F}_2$ & $\mathcal{H}_5\,[31,26,3]$ & Dot & 21 & 21 & $[52,26]$ & 8 & 10--12 \\
$\mathbb{F}_3$ & BCH $[13,7,5]$ & Dot & 7 & 7 & $[20,7]$ & 9 & 9 \\
$\mathbb{F}_4$ & $[10,5]$ & Hermitian & 5 & 5 & $[15,5]$ & 6 & 8 \\
$\mathbb{F}_5$ & $[12,6]$ & Dot & 6 & 6 & $[18,6]$ & 8 & 10 \\
\bottomrule
\end{tabular}%
}
\end{table}

The discovery of perfect partitions in $\Lambda_7(1^4)$ and $\Lambda_{15}(1^{12})$---extending Mollard's boundary---shows that the system can handle irregular combinatorial structures. Table~\ref{tab:lucas-summary} catalogs the generalized Lucas cubes tested, detailing the number of centers and the varying ball radii required for the perfect partitions. The table highlights that SolEvolve successfully found SAT instances for highly irregular topologies, such as $\Lambda_{15}(1^{12})$ and $\Lambda_{15}(1^{11})$, proving the existence of perfect partitions with non-uniform ball sizes (radii 14, 15, and 16). Conversely, it provided rigorous UNSAT proofs for smaller forbidden sequences like $s=4$, demonstrating a sharp transition in the covering properties of these topological subgraphs.

\begin{table}[t]
\centering
\caption{Perfect partition results for generalized Lucas cubes}
\label{tab:lucas-summary}
\resizebox{\columnwidth}{!}{%
\begin{tabular}{llllll}
\toprule
$(n, s)$ & $|V|$ & Ball Sizes & Centers & Status & Time \\
\midrule
$(7, 4)$ & 99 & $\{6, 7, 8\}$ & 15 & SAT & $<0.1$s \\
$(7, 3)$ & 71 & --- & --- & UNSAT & 0.02s \\
$(7, 2)$ & 29 & --- & --- & UNSAT & $<0.01$s \\
$(15, 12)$ & 32,707 & $\{14, 15, 16\}$ & 2,047 & SAT & 5.6s \\
$(15, 11)$ & 32,647 & $\{14, 15, 16\}$ & 2,047 & SAT & 388s \\
$(15, 10)$ & 32,527 & $\{14, 15, 16\}$ & $\le 2047$ & Timeout & $>$2hr \\
$(15, 4)$ & 18,842 & --- & --- & UNSAT & 496s \\
$(15, 3)$ & 9,327 & --- & --- & UNSAT & 5.4s \\
$(15, 2)$ & 1,364 & --- & --- & UNSAT & $<0.1$s \\
\bottomrule
\end{tabular}%
}
\end{table}

\subsection{Ablation study: component contributions}\label{sec:ablation}

We evaluate the contribution of each SolEvolve component through controlled experiments across multiple domains. Table~\ref{tab:ablation} summarizes the key findings.

\begin{table}[t]
\centering
\small
\caption{Ablation study across domains. For $[22,11,7]$: success rate over 50 runs. For $[43,10,16]$: $A_d$ reduction from BKLC ($A_{16}=225$). Weight variance measures structural uniformity.}
\label{tab:ablation}
\begin{tabular}{@{}lcc@{}}
\toprule
Component Configuration & Metric & Result \\
\midrule
\multicolumn{3}{l}{\textit{Binary Code: $[22,11,7]$ SAT-GA Repair (50 runs)}} \\
\quad Standard GA (fitness only) & Success & 2\% \\
\quad GA + SAT Repair (no Reflector) & Success & 4\% \\
\quad Full SolEvolve & Success & \textbf{10\%} \\
\quad Full SolEvolve & Wt. var. & $5.44 \pm 0.82$ \\
\midrule
\multicolumn{3}{l}{\textit{$A_d$ Optimization: $[43,10,16]$}} \\
\quad Single SAT query & Best $A_{16}$ & 91 \\
\quad Iterative blocking (no Reflector) & Best $A_{16}$ & 84 \\
\quad Full SolEvolve (6 iter.) & Best $A_{16}$ & \textbf{80} (64\%$\downarrow$) \\
\midrule
\multicolumn{3}{l}{\textit{Self-Orthogonal: Hamming $[15,11,3]$ Embedding}} \\
\quad Rank-one elimination only & $d_{\min}$ & 4 \\
\quad Full SolEvolve (SAT + $d$ constr.) & $d_{\min}$ & \textbf{6} \\
\bottomrule
\end{tabular}
\end{table}

\paragraph{Hybrid SAT-GA repair.} The binary $[22,11,7]$ code experiment validates the Reflector's semantic feedback. Across 50 runs, full SolEvolve achieved 10\% repair success---upgrading $[22,11,6]$ candidates to $[22,11,7]$---versus 2\% for fitness-only GA. Successful repairs exhibited weight variance of $5.44 \pm 0.82$, significantly lower than the standard code's variance of 21.0, indicating the Reflector optimizes structural uniformity beyond distance satisfaction.

\paragraph{Iterative $A_d$ optimization.} For $[43,10,16]$, a single SAT query achieved $A_{16}=91$; iterative blocking without Reflector guidance reached $A_{16}=84$. The full framework, with the Reflector analyzing weight patterns and proposing symmetry-breaking constraints, achieved $A_{16}=80$---a 64\% reduction from BKLC's $A_{16}=225$. All six codes were verified inequivalent via Magma.

\paragraph{Algorithm vs.\ constraint optimization.} Self-orthogonal embedding reveals complementary roles: the Generator's rank-one elimination achieves optimal length ($n + \rank(GG^T)$) but yields $d=4$; adding SAT-based distance constraints achieves $d=6$, matching the shortened Golay code.

\subsection{Addressing data contamination concerns}\label{sec:contamination}

A critical question for LLM-derived algorithms is whether the Generator genuinely \textit{conceived} the rank-one elimination algorithm or merely \textit{retrieved} it from training data. While Kim et al.~\cite{Kim2021} and related works on self-orthogonal embeddings exist in the literature, we provide evidence for genuine synthesis:

\paragraph{Synthetic problem test.} We constructed a ``twisted embedding'' problem with no prior literature: given a code $\mathcal{C}$ over $\mathbb{F}_4$ with generator matrix $G$, find the shortest code $\tilde{\mathcal{C}}$ such that $G\bar{G}^T + \omega \cdot SS^T = \mathbf{0}$ (where $\bar{G}$ denotes Hermitian conjugation and $\omega$ is a primitive element). This twisted self-orthogonality condition has no documented solution. The Generator, given only the problem definition, derived a correct algorithm that reduces to rank-one elimination over $\mathbb{F}_4$ with modified update rules. The solution was verified correct via Magma.

\paragraph{Novelty of formulation.} The rank-one elimination algorithm as presented (Algorithm~\ref{alg:rank1}) differs from Kim et al.'s self-orthogonality matrix approach~\cite{Kim2022} in its iterative Gram-matrix elimination structure. While the theoretical lower bound ($n + \rank(GG^T)$) was known, the specific algorithmic realization via anisotropic/hyperbolic case analysis represents a novel synthesis.

\paragraph{Prompt ablation.} When we removed mathematical context from the Generator's prompt (providing only input/output specifications without mentioning Gram matrices or quadratic forms), the Generator still produced a functionally equivalent algorithm, suggesting it derived the structure from the problem constraints rather than pattern-matching to known solutions.

\subsection{Limitations and future work}

While SolEvolve demonstrates measurable improvements across multiple combinatorial domains, several limitations warrant discussion.

\paragraph{Scalability constraints.} The framework's computational requirements grow substantially with problem size. For binary linear codes, instances with $n > 50$ become challenging for single-threaded SAT solving, as evidenced by the solve time increase from 5.6 seconds for $\Lambda_{15}(1^{12})$ to 388 seconds for $\Lambda_{15}(1^{11})$. We note that $n=50$ is relatively small in coding theory; extending to truly open problems (where $n$ is larger) remains a challenge, highlighting the need for parallel solving strategies.

\paragraph{Threshold sensitivity.} The Reflector's state transitions depend on thresholds ($\tau_\rho = 0.01$, $\tau_D = 0.15$, $N_{\text{stag}} = 50$) determined through preliminary experiments. We conducted sensitivity analysis: varying $\tau_\rho \in [0.005, 0.02]$ and $N_{\text{stag}} \in [30, 70]$ changed the binary $[22,11,7]$ code success rate by $\pm 2\%$, suggesting moderate robustness. However, the diversity threshold $\tau_D$ showed higher sensitivity ($\pm 4\%$ when varied in $[0.10, 0.20]$), indicating that population diversity monitoring requires careful calibration.

\paragraph{Diversity evolution.} The diversity index $D_t$ typically follows a characteristic pattern: initial high diversity ($D_0 \approx 0.4$), convergence during exploitation ($D_t \to 0.15$ by generation 30--40), and recovery after SAT repair intervention ($D_t \to 0.25$). We did not observe complete diversity collapse ($D_t < 0.05$) in successful runs, suggesting the Reflector's interventions effectively prevent premature convergence. Full $D_t$ trajectories are provided in the supplementary material.

\paragraph{LLM component limitations.} The Generator's reasoning process remains partially opaque. While the rank-one elimination algorithm emerged through LLM-driven exploration, we cannot fully explain the intermediate reasoning steps. The contamination study (Section~\ref{sec:contamination}) provides evidence for genuine synthesis, but we cannot definitively rule out retrieval of related concepts recombined in novel ways. Additionally, the Reflector occasionally misinterprets Verifier feedback---for instance, confusing numerical stagnation with fundamental infeasibility---leading to premature termination in approximately 5\% of runs.

\paragraph{Algorithm conception vs. distance optimization.} The rank-one elimination algorithm, while provably optimal for embedding length, cannot directly manage the minimum distance of the resulting code. As demonstrated in Section~\ref{ch:so}, the algorithm yields $d=4$ for binary Hamming embeddings, whereas SAT-based approaches achieve $d=6$ by incorporating distance constraints. This limitation necessitates a multi-component solution within the current framework: the Generator conceives the base algorithm optimizing the primary objective (embedding length), while the Evolver and SAT-based modules refine solutions for secondary objectives (minimum distance, weight distribution).

\paragraph{Multi-objective considerations.} The tension between embedding length and minimum distance exemplifies a broader challenge in automated algorithm discovery: balancing competing objectives. Current LLM-based approaches, including SolEvolve, typically optimize single objectives sequentially rather than simultaneously. Integrating multi-objective optimization criteria---inspired by works such as ``Large Language Model for Multi-objective Evolutionary Optimization''~\cite{Liu2023LLMMOO}---represents a promising direction for future development. Specifically, we envision extending the Generator's prompting structure to include Pareto-aware trade-off reasoning, enabling the system to propose solutions that explicitly balance length minimality against distance maximization. This would require augmenting the Reflector's feedback vocabulary to communicate multi-objective trade-offs (e.g., ``embedding achieves $n+7$ but $d=4$; consider sacrificing 2 columns for potential $d=6$'').

\paragraph{Reproducibility and external dependencies.} The framework's reliance on external databases (e.g., \texttt{codetables.de}) for bound verification introduces fragility: network failures or database updates may affect result interpretation. We mitigate this through caching and fallback mechanisms, but a fully self-contained verification module remains a goal for future development.

\paragraph{Future directions.} Several promising extensions emerge from this work. First, the integration of portfolio SAT solving~\cite{Zhang2025SolSearch} could accelerate convergence on difficult instances. Second, extending the Generator's capabilities to conceive algorithms for non-binary fields and ring structures would broaden the framework's mathematical scope. Third, developing interpretability tools for the Reflector's semantic analysis would enhance trust in autonomous discovery claims. Finally, investigating transfer learning across combinatorial domains---where insights from code discovery inform MOLS construction, for example---could unlock synergies currently unexploited by domain-specific optimization.
